\begin{ExerciseBlock}
\ExerciseHeading{05}{Federkonstante aus statischer Einzelmessung}{laboratory_measurements_static_spring_constant}
Eine Feder wird an einem Stativ aufgehängt. Vor und nach Anhängen einer Masse $m$ werden die Höhen $h_1$ und $h_2$ des Federendes über der Tischplatte gemessen. Aus $\Delta L=h_1-h_2$ folgt die Federkonstante $D=mg/\Delta L$. Verwenden Sie die lineare Abschätzung der absoluten Unsicherheit.
\begin{enumerate}
\item Gemessen wurden $h_1=\qty{345}{\milli\metre}$ und $h_2=\qty{281}{\milli\metre}$. Setzen Sie für beide Messungen $u(h)=\qty{1}{\milli\metre}$ an und geben Sie die Höhen mit ihrer absoluten Unsicherheit an.
\item Geben Sie die Verlängerung $\Delta L$ mit ihrer absoluten Unsicherheit an.
\item Recherchieren Sie die Werte der Erdbeschleunigung am Äquator und an den Polen. Geben Sie einen Wert für $g$ mit einer absoluten Unsicherheit an, die diesen Bereich abdeckt.
\item Eine Laborwaage zeigt $m=\qty{200}{\gram}$ ohne Nachkommastellen an. Geben Sie $m$ mit einer anhand der Anzeige abgeschätzten Unsicherheit an und nennen Sie Ihre Annahme.
\item Berechnen Sie $D$ und die absolute Unsicherheit $u(D)$.
\item Runden Sie Unsicherheit und Ergebnis auf zueinander passende Stellen.
\end{enumerate}
\end{ExerciseBlock}
\UseExerciseSolution{laboratory_measurements_static_spring_constant}
